% !PRINTED 34
\weberpage{34}
%<<CONT
\weq{3}{1 = A_1^{\alpha_1' - \alpha_1} A_2^{\alpha_2' - \alpha_2} \ldots A_n^{\alpha_n' - \alpha_n},}
also eine Darstellung des Elementes 1. Also können wir nach (2) setzen:
\[
\alpha_1' = \alpha_1 + h_1,\ \alpha_2' = \alpha_2 + h_2 \ldots \alpha_n' = \alpha_n + h_n.
\]
Es kann folglich auch nicht mehr verschiedene Darstellungen des Elementes $\Theta$ geben, als die Anzahl der Darstellungen des Elementes 1 beträgt. Wir bezeichnen die Anzahl dieser Darstellungen mit $h$. Im Ganzen ist aber die Anzahl der verschiedenen möglichen Bestimmungen der $\alpha$ in (1) gleich dem Producte $a_1 a_2 \ldots a_n$, und da $n$ die Anzahl der Elemente von $S$ ist, so folgt
\weq{4}{n h = a_1 a_2 \ldots a_n.}
Aus der Formel (4) ergiebt sich der Satz:

\gesperrt{2. Wenn $r$ eine im Grade $n$ der Gruppe aufgehende Primzahl ist, so giebt es in $S$ ein Element vom Grade $r$.}

Denn aus (4) sieht man zunächst, dass einer der Grade $a_1, a_2, \ldots a_n$ durch $r$ theilbar ist. Ist also etwa $a_1 = rk$, so ist $A_1^k$ ein Element vom Grade $r$.

\gesperrt{3. Sind $A, B \ldots$ irgend welche Elemente in $S$ von den Graden $a, b \ldots$, so ist auch $AB \ldots$ ein Element in $S$, und der Grad dieses Productes ist ein Theiler des kleinsten gemeinschaftlichen Vielfachen $m$ von $a, b \ldots$}

Denn es ist
\[
(AB \ldots)^m = A^m B^m \ldots = 1,
\]
und folglich $m$ ein Vielfaches des Grades von $AB \ldots$

\gesperrt{4. Ist der Grad $n$ der Gruppe $S$ in zwei Factoren $n = ab$ zerlegt, so dass $a$ und $b$ relativ prim sind, so giebt es in $S$ genau $a$ Elemente $A$, deren Grad ein Theiler von $a$ ist, und $b$ Elemente $B$, deren Grad ein Theiler von $b$ ist, und in der Form}
\weq{5}{\Theta = AB}
\gesperrt{sind sämmtliche Elemente von $S$ und jedes nur einmal enthalten.}

Um die Richtigkeit dieses Satzes einzusehen, stellen wir folgende Ueberlegung an. Der Inbegriff $\mathfrak{A}$ der Elemente $A$, deren
%CONT>>
